\section{Supplementary Figures}
\label{app:figures}

\begin{figure}[htbp]
  \centering
  \includegraphics[width=\linewidth]{SOM_swiss}
  \caption{PCA-coloured landscape patch map of Switzerland.
  Codebook vectors are projected onto 3 principal components
  (mapped to RGB), so perceptually similar colours indicate
  environmentally similar patches. Major biogeographic zones
  are clearly distinguishable: lowland plateaus (Swiss Plateau,
  Po Valley) share warm tones, high-alpine areas appear as a
  distinct cool cluster, and forested pre-alpine slopes form
  intermediate gradients.}
  \label{fig:pca-patches}
\end{figure}

\begin{figure}[htbp]
  \centering
  \includegraphics[width=0.8\linewidth]{patch_size_distribution}
  \caption{Patch size distribution (log scale) for the 166\,464
  landscape patches. The vertical dashed line marks the minimum
  threshold of 100~pixels (${\sim}0.09$\,km$^2$).}
  \label{fig:patch-stats}
\end{figure}

% Additional supplementary figures below this line

\section{Random Forest Baseline Details}
\label{app:rf-baseline}

\Cref{tab:rf-summary} summarises the cross-validated performance
metrics for the Random Forest baseline across all 2,696 species.

\begin{table}[htbp]
  \centering
  \caption{Random Forest baseline: summary statistics across
  2,696 species (5-fold cross-validation, threshold\,=\,0.5).}
  \label{tab:rf-summary}
  \begin{tabular}{lrrr}
    \toprule
    Metric & Mean & Median & Std \\
    \midrule
    AUC       & 0.874 & 0.877 & 0.056 \\
    Accuracy  & 0.832 & 0.831 & 0.042 \\
    Precision & 0.733 & 0.754 & 0.151 \\
    Recall    & 0.481 & 0.486 & 0.185 \\
    F1        & 0.561 & 0.585 & 0.173 \\
    MCC       & 0.498 & 0.512 & 0.163 \\
    TSS       & 0.430 & 0.433 & 0.176 \\
    \bottomrule
  \end{tabular}
\end{table}

The AUC distribution is strongly right-concentrated: 99.3\,\% of
species achieve AUC\,$>$\,0.7, and 91.6\,\% exceed 0.8.

\begin{figure}[htbp]
  \centering
  \includegraphics[width=\linewidth]{codebook_epoch_150}
  \medskip
  \includegraphics[width=\linewidth]{codebook_epoch_200}
  \caption{SOM codebook feature maps at epoch~150 (top) and
  epoch~200 (bottom). At 150~epochs the codebook shows smooth
  gradients with distinct regional structure across all 12
  features. At 200~epochs, overtraining causes feature maps to
  fragment into noise---particularly visible in slope, TWI, and
  bio15---as the reduced neighbourhood radius allows individual
  neurons to chase local samples rather than preserving global
  topology.}
  \label{fig:codebook-comparison}
\end{figure}

\section{GNN-SDM Training Diagnostics}
\label{app:gnn-training}

\begin{figure}[htbp]
  \centering
  \includegraphics[width=0.85\linewidth]{gnn_training_val}
  \caption{Validation AUC curves during GNN-SDM training for
  the 12 selected species (6~common, 6~vulnerable). Early
  stopping restores the model at peak validation AUC. Most
  species converge within 200--350 epochs; common species with
  more training data tend to use the full 500-epoch budget.}
  \label{fig:gnn-val-auc}
\end{figure}

\begin{figure}[htbp]
  \centering
  \includegraphics[width=0.85\linewidth]{gnn_training}
  \caption{Training and validation loss (weighted MSE) for the
  12 selected species. The gap between training and validation
  loss remains small throughout, indicating minimal overfitting.
  Dropout ($p = 0.2$) and early stopping provide sufficient
  regularisation.}
  \label{fig:gnn-train-loss}
\end{figure}

\section{Per-Species Habitat Suitability Maps}
\label{app:suitability-maps}

\begin{figure}[htbp]
  \centering
  \includegraphics[width=\linewidth]{suitablity_per_species}
  \caption{GNN-SDM habitat suitability maps for the 12 selected
  species. Warm colours indicate high predicted suitability.
  Common species (top two rows) show broader predicted ranges
  consistent with their generalist niches, while vulnerable
  species (bottom two rows) display concentrated suitability
  in restricted habitat patches.}
  \label{fig:suitability-per-species}
\end{figure}
